\section{Limitations}
\label{sec:b-limitations}

\begin{itemize}[leftmargin=1.4em]
\item \textbf{Two-checkpoint scope.} The controller acts at most once, across at
  most two online checkpoints, under a frozen continuation.  Deployment value
  describes a controller that meets the design's checkpoint rule; a stopping-time
  or every-checkpoint deployment adds selection, coverage and post-intervention
  support terms the design does not estimate.
\item \textbf{Replay allowance is an assumption.} The $2p$ population shift is a
  sensitivity assumption on the registered grid, never a measured quantity
  (\cref{sec:b-theory}).  An understated $p$ is a sensitivity failure, not a
  refutation of the conditional proposition; zero misses in a finite audit panel
  do not certify a small $p$.
\item \textbf{Task mechanism.} The candidate-consistency mechanism classes
  establish that regime differences can exist under specified conditions; they do
  not establish prevalence in natural deployments.  The screen law is an
  artificial diagnostic law, not natural prevalence.
\item \textbf{Feasibility is a registry claim.} The mask is only as valid as its
  registry and evaluator.  Withholding verifier access synthetically is a
  capacity condition, not evidence of natural unavailability.
\item \textbf{Frozen continuation and representation.} Values are relative to the
  registered $\rho$ and horizon.  The representation may discard decision-relevant
  history; representation adequacy is separate and not assumed.
\item \textbf{Finite data and multiplicity.} Effective sample sizes at small
  capacities can be much smaller than root counts.  A cluster bootstrap is
  asymptotic and is not a distribution-free finite-sample certificate; no
  per-root i.i.d.\ certificate is claimed.
\item \textbf{Substrate and external validity.} Effects are scoped to the
  specified model, harness, continuation, verifier and task distribution.  No
  effect for another model, prompt, tool policy or workload is identified.
\item \textbf{Safety language.} Empirical harm endpoints, certified
  noninferiority and individual harm probability are distinct claims with
  distinct evidence requirements; none is upgraded by a passing test or a LaTeX
  build.
\item \textbf{Batch versus streaming.} Batch top-$q$ results do not establish
  streaming rate guarantees; the two are reported separately.
\end{itemize}
